\documentclass{article}
\usepackage{graphicx}
\usepackage{amsmath}

\title{Autocorrelation Analysis in Florida weather}
\author{Bridget Smith, Anna Cavalieri Canosa, Lehan Geng, Zhilu Zhang}
\date{November 2024}

\begin{document}

\maketitle

 \section{Introduction}
 \vspace{-0.5em}
This document presents the results of the autocorrelation analysis of Key West annual mean temperatures. The goal of this analysis was to determine whether the temperatures of one year are significantly correlated with those of the following year across the entire time series.

\section{Methods of Analysis}
\vspace{-0.5em}
The analysis consists of two parts. First, the original correlation coefficient between successive years was calculated. Second, the statistical significance of the correlation was assessed through 10,000 simulations. The frequency of randomly permuted correlations exceeding the original correlation was recorded to evaluate its significance.
\section{Results \& Conclusion}
\vspace{-0.5em}
{\selectfont\subsubsection*{Results}
\vspace{-0.5em}
The original correlation coefficient between successive years' temperatures was found to be 0.326. In the 10,000 simulations, only 6 of the random correlation coefficients were greater than the original correlation. The approximate p-value is therefore:
\vspace{-0.5em}
\[
\text{p-value} = \frac{6}{10000} = 0.0006
\]}

{\selectfont\subsubsection*{Interpretation}
The original correlation coefficient of 0.326169 indicates a positive but weak linear relationship between consecutive years.
The p-value of 0.0006 is much smaller than the commonly used significance level of 0.05, indicates that the original correlation is statistically significant and unlikely to have occurred by chance.}

{\selectfont\subsubsection*{Conclusion}
The results of this analysis demonstrates that temperatures of one year are significantly correlated with those of the next year in the given location.}

\end{document}
 
